\blankpage

\thispagestyle{empty}
\clearpage
\begin{tikzpicture}[remember picture, overlay]
\node at (3,-7) {\includegraphics[width=8cm]{./recettes/edités/Biscuit vanille et Biscuit à l'anis.jpg}}; 
\end{tikzpicture}\clearpage\thispagestyle{empty}
% \begin{tikzpicture}[remember picture,overlay]
% \fill[black] (-3,3.5) rectangle (14,-21.5);
% \end{tikzpicture}
\clearpage

\chapter*{Biscuits à l’anis}

\section{Ingrédients originaux}

\begin{itemize}
\item 3 œufs
\item 1/2 verre huile
\item 150 gr. de sucre
\item 1 poignée d’anis
\item Farine autant que ça prend
\end{itemize}

\section{Préparation adaptée}

Battez les œufs avec le sucre jusqu’à obtenir un mélange clair et homogène, ajoutez l’huile puis l’anis. Incorporez la farine progressivement jusqu’à obtenir une pâte molle, souple et légèrement collante. Ajoutez la levure chimique (\textsc{b.\,p.}) en comptant environ 1 cuillère à café pour chaque verre de farine utilisé. Formez des petites boules ou des biscuits légèrement aplatis, disposez-les sur une plaque et faites cuire au four préchauffé à 180°\,\textsc{c} pendant 15 à 25 minutes, jusqu’à ce qu’ils soient légèrement dorés.

\chapter{Biscoitos de anis}

\section{Ingredientes originais}

\begin{itemize}
\item 3 ovos
\item 1/2 copo de óleo
\item 150 g de açúcar
\item 1 punhado de anis
\item Farinha quanto baste
\end{itemize}

\section{Modo de preparo adaptado}

Bata os ovos com o açúcar até obter uma mistura clara e homogênea, adicione o óleo e depois o anis. Incorpore a farinha aos poucos até obter uma massa macia, maleável e levemente pegajosa. Adicione o fermento químico calculando cerca de 1 colher de chá para cada copo de farinha utilizado. Modele pequenas bolinhas ou biscoitos levemente achatados, disponha-os em uma assadeira e leve ao forno preaquecido a 180°,\textsc{c} por 15 a 25 minutos, até ficarem levemente dourados.